
\section{Xây dựng mô hình Random Forest}
- feature engineering: 
Đầu tiên tôi phân tích absolute correlation
Lý do: bỏ đi một cái 
*_minutes và *_charge (biến số)Dùng tương quan Pearson (hoặc Spearman).
\(H_0\): \(\rho = 0\), \(H_1\): \(\rho \neq 0\).
Vì \(|r| \approx 1\), hai biến gần như trùng thông tin, giữ minutes, bỏ charge để tránh dư thừa.

-> chắc dựa vào bảng rồi viết thôi

Top absolute correlations:
            feature_1          feature_2  correlation  abs_correlation
    total_day_minutes   total_day_charge     1.000000         1.000000
    total_eve_minutes   total_eve_charge     1.000000         1.000000
  total_night_minutes total_night_charge     0.999999         0.999999
   total_intl_minutes  total_intl_charge     0.999993         0.999993
       account_length    total_day_calls     0.028240         0.028240
  total_night_minutes  total_night_calls     0.026972         0.026972
    total_night_calls total_night_charge     0.026949         0.026949
number_vmail_messages   total_eve_charge     0.019496         0.019496
number_vmail_messages  total_eve_minutes     0.019490         0.019490
     total_day_charge total_intl_minutes    -0.019490         0.019490
    total_day_minutes total_intl_minutes    -0.019486         0.019486
     total_day_charge  total_intl_charge    -0.019419         0.019419

- thứ hai tôi đi phân tích Target association
Numeric correlation with churn=yes:
                      feature  correlation_with_churn_yes  abs_correlation
number_customer_service_calls                    0.212564         0.212564
            total_day_minutes                    0.207705         0.207705
             total_day_charge                    0.207700         0.207700
        number_vmail_messages                   -0.097633         0.097633
            total_eve_minutes                    0.089288         0.089288
             total_eve_charge                    0.089282         0.089282
           total_intl_minutes                    0.063285         0.063285
            total_intl_charge                    0.063275         0.063275
             total_intl_calls                   -0.046893         0.046893
          total_night_minutes                    0.045677         0.045677
           total_night_charge                    0.045673         0.045673
               account_length                    0.021203         0.021203
              total_day_calls                    0.016130         0.016130
            total_night_calls                   -0.006986         0.006986
              total_eve_calls                   -0.006284         0.006284

Highest categorical churn-rate groups with at least 20 records:
           feature      category  total  churn_yes  churn_rate
         area_code area_code_510   1246        184      0.1477
         area_code area_code_408   1259        177      0.1406
         area_code area_code_415   2495        346      0.1387
international_plan           yes    473        199      0.4207
international_plan            no   4527        508      0.1122
             state            CA     52         14      0.2692
             state            NJ    112         28      0.2500
             state            WA     98         24      0.2449
             state            TX    116         26      0.2241
             state            MT     99         21      0.2121
             state            MD    102         21      0.2059
             state            NV     90         17      0.1889
             state            ME    103         19      0.1845
             state            KS     99         18      0.1818
             state            OK     90         16      0.1778

=> international_plan là tín hiệu mạnh nhất ở nhóm categorical
Churn của yes = 42.07% vs no = 11.22%.
Đây là phân khúc rủi ro rất cao, nên ưu tiên retention.
Với numeric, mạnh nhất là hành vi “không hài lòng + dùng nhiều”
number_customer_service_calls corr = 0.2126 (cao nhất).
total_day_minutes corr = 0.2077 (và total_day_charge tương đương, nhưng là cột dư thừa).
Hàm ý: khách gọi CSKH nhiều và dùng ban ngày cao dễ churn hơn.
number_vmail_messages tương quan âm (-0.0976)
Khách dùng voicemail nhiều có xu hướng ít churn hơn (mức vừa phải).
area_code gần như không tạo khác biệt thực tế
0.1387–0.1477, rất sát mức nền churn (~14.14%).
Không nên kỳ vọng area_code là feature mạnh.
state có vài bang churn cao (CA, NJ, WA, TX...) nhưng cần thận trọng
Một số bang có mẫu nhỏ (ví dụ CA chỉ 52), dễ nhiễu.
Dùng để theo dõi thị trường, không nên quyết định lớn chỉ dựa vào state đơn lẻ.

Giữ/ưu tiên: number_customer_service_calls, international_plan, total_day_minutes, number_vmail_messages.

- Total missing cells: 0
Duplicate rows: 0

Top outlier counts:
                      feature      q1     q3    iqr  lower_bound  upper_bound  outlier_count  outlier_percent
number_customer_service_calls   1.000   2.00  1.000      -0.5000       3.5000            399             7.98
             total_intl_calls   3.000   6.00  3.000      -1.5000      10.5000            118             2.36
           total_intl_minutes   8.500  12.00  3.500       3.2500      17.2500             72             1.44
            total_intl_charge   2.300   3.24  0.940       0.8900       4.6500             72             1.44
        number_vmail_messages   0.000  17.00 17.000     -25.5000      42.5000             60             1.20
            total_eve_minutes 166.375 234.10 67.725      64.7875     335.6875             43             0.86
            total_night_calls  87.000 113.00 26.000      48.0000     152.0000             43             0.86
             total_eve_charge  14.140  19.90  5.760       5.5000      28.5400             42             0.84
          total_night_minutes 166.900 234.70 67.800      65.2000     336.4000             39             0.78
           total_night_charge   7.510  10.56  3.050       2.9350      15.1350             39             0.78

           Chất lượng dữ liệu tốt: Total missing cells = 0, Duplicate rows = 0.
Bảng outlier dùng quy tắc IQR:
q1, q3: tứ phân vị 25% và 75%
iqr = q3 - q1
lower_bound = q1 - 1.5*iqr, upper_bound = q3 + 1.5*iqr
outlier_count/percent: số và tỷ lệ điểm vượt ngưỡng
Insight chính:

number_customer_service_calls có outlier cao nhất (7.98%) và trước đó cũng tương quan mạnh với churn -> nhiều khả năng là tín hiệu hành vi thật, không phải lỗi dữ liệu.
Nhóm total_intl_* có outlier vừa phải (1–2%) -> có thể là nhóm khách dùng quốc tế nhiều.
*_charge và *_minutes cùng xuất hiện vì hai cột này gần như trùng thông tin (đã thấy correlation ~1), nên giữ một bên là đủ.


- Top numeric tests by Mann-Whitney p-value:
                      feature    mean_no   mean_yes  mean_difference_yes_minus_no  cohens_d  mannwhitney_pvalue
            total_day_minutes 175.746564 207.870580                     32.124016  0.536657            0.000000
             total_day_charge  29.877494  35.338416                      5.460922  0.536637            0.000000
number_customer_service_calls   1.457722   2.254597                      0.796875  0.522435            0.000000
        number_vmail_messages   8.291870   4.496464                     -3.795407 -0.300835            0.000000
            total_eve_minutes 198.805031 211.757850                     12.952819  0.254649            0.000000
             total_eve_charge  16.898654  17.999562                      1.100908  0.254633            0.000000
             total_intl_calls   4.481947   4.151344                     -0.330604 -0.132520            0.000004
           total_intl_minutes  10.190869  10.692362                      0.501493  0.180935            0.000049
            total_intl_charge   2.752055   2.887426                      0.135371  0.180912            0.000049
           total_night_charge   8.975593   9.273607                      0.298014  0.132054            0.001668

Categorical chi-square tests:
           feature       chi2   pvalue  dof  cramers_v  min_expected_count
international_plan 333.186449 0.000000    1   0.258142             66.8822
   voice_mail_plan  60.552418 0.000000    1   0.110048            187.0722
             state  96.899041 0.000079   50   0.139211              7.3528
         area_code   0.562981 0.754658    2   0.010611            176.1844
Numeric tests (Mann-Whitney)
pvalue rất nhỏ nghĩa là nhóm churn yes và no khác nhau có ý nghĩa thống kê ở nhiều biến.
Nhưng quan trọng hơn là cohens_d (độ lớn ảnh hưởng):
Mạnh nhất và mức vừa: total_day_minutes (d≈0.54), number_customer_service_calls (d≈0.52).
Vừa/nhỏ: number_vmail_messages (d≈-0.30, dấu âm = churn có voicemail ít hơn), total_eve_minutes (d≈0.25).
Nhỏ: các biến intl/night (d ~0.13–0.18).
Categorical tests (Chi-square)
international_plan: rất mạnh (p<0.001, Cramer's V≈0.258) -> feature categorical quan trọng nhất.
voice_mail_plan: có ý nghĩa nhưng yếu hơn (V≈0.11).
state: có ý nghĩa thống kê nhưng hiệu ứng yếu-vừa (V≈0.139), dùng được nhưng không phải tín hiệu chính.
area_code: không ý nghĩa (p=0.755, V≈0.011) -> gần như không hữu ích.
Insight hành động
Ưu tiên feature: international_plan, number_customer_service_calls, total_day_minutes, number_vmail_messages.
Bỏ các cột *_charge nếu đã giữ *_minutes (vì trùng thông tin gần hoàn hảo).
area_code có thể loại khỏi mô hình baseline.

- sau khi bỏ đi 4 biến *charge thì tôi vẫn còn 16 biến. Tôi sẽ tiến hành lập model RF baseline với 16 biến.

- ## 2. Với biến nhị phân thì không nên giữ cả yes và no

Bạn đang có:

text
international_plan_no
international_plan_yes
voice_mail_plan_no
voice_mail_plan_yes


Thực ra chỉ cần giữ một cột là đủ, ví dụ:

text
international_plan_yes
voice_mail_plan_yes


Vì:

text
international_plan_no = 1 - international_plan_yes
voice_mail_plan_no = 1 - voice_mail_plan_yes


Giữ cả hai không sai nghiêm trọng với Random Forest, nhưng bị dư thừa.

Nên dùng:

text
international_plan_yes
voice_mail_plan_yes


- với biến state: 
“state does not improve predictive performance under our RF setup; removing it yields higher holdout/CV F1.”

- state và area \_code: bỏ để giảm nhiễu/overfit và tăng khả năng tổng quát
area \_code hầu như không liên hệ churn: p = 0.7547, Cramer's V = 0.0106 (rất yếu).
state có p nhỏ nhưng hiệu ứng thực tế yếu hơn (V = 0.139) và là biến định danh nhiều mức (mỗi state chỉ khoảng 52–158 mẫu), dễ làm cây tách theo nhiễu cục bộ.

=> (1) thống kê độc lập giữa area_code và churn, (2) ablation cho thấy thêm area_code/state không cải thiện F1-CI.

Cách giải thích thuyết phục cho area_code:
Kiểm định liên hệ với churn\(H_0\): area_code độc lập với churn.
Kết quả: p = 0.7547 (> 0.05) ⇒ không bác bỏ \(H_0\), tức là chưa có bằng chứng area_code liên hệ churn.

Độ lớn hiệu ứngCramer's V = 0.0106 ⇒ gần 0, mức rất yếu/không đáng kể.

Mô tả trực quan
Tỷ lệ churn theo 3 nhóm area code gần như giống nhau (~13.9%\–14.8\%), chênh lệch nhỏ.
Kết luận viết báo cáo:
“area_code không có ý nghĩa thống kê và hiệu ứng thực tế gần như bằng 0; giữ biến này chỉ tăng số chiều, dễ thêm nhiễu cho RF nên loại bỏ để mô hình gọn và tổng quát hơn.”
Cho state thì giải thích tương tự nhưng thêm ý:
Dù có thể p nhỏ, hiệu ứng thực tế yếu (Cramer's V nhỏ), số mức nhiều (50+) nên dễ overfit theo nhóm nhỏ, khó tổng quát và khó hành động nghiệp vụ.
=> Loại state là quyết định theo nguyên tắc parsimony + generalization.

international \_plan \_yes, voice \_mail \_plan \_yes mã hóa 0/1 là chuẩn
Đây là biến nhị phân rõ nghĩa nghiệp vụ.
Mã hóa yes=1, no=0 giúp mô hình đọc trực tiếp.
Bỏ cột * \_no để tránh trùng thông tin hoàn toàn với * \_yes.


=> cách chứng minh: area_code vs churn (biến phân loại)Kiểm định: Chi-square independence.
\(H_0\): area_code độc lập với churn.
\(H_1\): area_code có liên hệ với churn.
Mức ý nghĩa: \(\alpha = 0.05\).
Kết quả của bạn: p = 0.7547, Cramer's V = 0.0106.
Kết luận: không bác bỏ \(H_0\), hiệu ứng rất nhỏ, loại area_code.

state vs churn (biến phân loại nhiều mức)Kiểm định: Chi-square independence.
\(H_0\): state độc lập với churn.
\(H_1\): state có liên hệ với churn.
Báo cáo thêm Cramer's V để đo độ mạnh hiệu ứng.
Dù có thể p nhỏ, nếu V nhỏ thì ý nghĩa thực tiễn thấp; thêm vào đó state nhiều mức dễ overfit, nên loại để tăng khả năng tổng quát.

- insight tổng
5000 dòng, churn=yes là 14.14\% (mất cân bằng lớp).
Không có missing/duplicate.
Các cột $* \_charge$ gần như trùng thông tin với $* \_minutes$ nên được khuyến nghị bỏ.
Tín hiệu churn nổi bật: number \_customer \_service \_calls, total \_day \_minutes, international \_plan.

Giữ minutes, bỏ charge (hoặc ngược lại, nhưng phải nhất quán).
Thường dùng danh sách drop:
total \_day \_charge, total \_eve \_charge, total \_night \_charge, total \_intl \_charge.

=> giữ lại 13 features: account_length
number_vmail_messages
total_day_minutes
total_day_calls
total_eve_minutes
total_eve_calls
total_night_minutes
total_night_calls
total_intl_minutes
total_intl_calls
number_customer_service_calls
international_plan_yes
voice_mail_plan_yes

Ablation cho thấy v2 (13 features) tốt hơn v1 (64 features)
max \_depth=5:
v1 CV F1 mean = 0.0209 vs v2 = 0.4720
v1 Holdout F1 = 0.0946 vs v2 = 0.4432
max \_depth=None:
v1 CV F1 mean = 0.7187 vs v2 = 0.8029
v1 Holdout F1 = 0.7647 vs v2 = 0.7869
--------------------------------------------------
Baseline modeling complete (max_depth=3).
Features used: 64
Train F1: 0.000000
Holdout F1: 0.000000
CV F1 mean (95% CI): 0.000000 [0.000000, 0.000000]
--------------------------------------------------
Baseline modeling complete (max_depth=5).
Features used: 64
Train F1: 0.081356
Holdout F1: 0.094595
CV F1 mean (95% CI): 0.020870 [-0.002786, 0.044525]
--------------------------------------------------
Baseline modeling complete (max_depth=None).
Features used: 64
Train F1: 1.000000
Holdout F1: 0.764706

--------------------------------------------------
Model v2 complete (max_depth=3).
Features used: 13
Train F1: 0.044828
Holdout F1: 0.107383
CV F1 mean (95% CI): 0.017392 [0.002253, 0.032531]
--------------------------------------------------
Model v2 complete (max_depth=5).
Features used: 13
Train F1: 0.641916
Holdout F1: 0.515464
CV F1 mean (95% CI): 0.479703 [0.430128, 0.529279]
--------------------------------------------------
Model v2 complete (max_depth=None).
Features used: 13
Train F1: 1.000000
Holdout F1: 0.824000
CV F1 mean (95% CI): 0.787214 [0.739351, 0.835078]

% *_minutes và *_charge (biến số)Dùng tương quan Pearson (hoặc Spearman).
% \(H_0\): \(\rho = 0\), \(H_1\): \(\rho \neq 0\).
% Vì \(|r| \approx 1\), hai biến gần như trùng thông tin, giữ minutes, bỏ charge để tránh dư thừa.
% Gợi ý dùng ngay cho mô hình:
% Với mô hình cây như Random Forest: không cần xóa outlier vội.
% Thử 2 phiên bản để so sánh F1:
% Giữ nguyên outlier
% Cap/winsorize (ví dụ chặn ở p99) cho vài cột lệch mạnh.
% Bỏ các cột * \_charge (dư thừa với * \_minutes).
% Giữ/ưu tiên: number \_customer \_service \_calls, international \_plan, total \_day \_minutes, number \_vmail \_messages.
% Thử feature tương tác: international \_plan × total \_intl \_minutes, customer \_service \_calls × total \_day \_minutes.
% Ưu tiên feature: international \_plan, number \_customer \_service \_calls, total \_day \_minutes, number \_vmail \_messages.
% Bỏ các cột * \_charge nếu đã giữ * \_minutes (vì trùng thông tin gần hoàn hảo).
% area \_code có thể loại khỏi mô hình baseline.





